%% ma: L10-B13-0069
%% lop: 10
%% chuong: 5
%% bai: 13
%% dang: 10.13.14
%% loai: TL
%% muc: VD
%% dap_an: "-"
%% nguon: "chuyen-de-cac-so-dac-trung-cua-mau-so-lieu-khong-ghep-nhom-toan-10-chuong-trinh-moi.pdf (GV. Phan Nhật Linh) (P03), Bài 02 Dạng 2 Bài tập tự luận 6"
%% kien_thuc: "Lớp 10 Bài 13 (số trung bình, trung vị, tứ phân vị, mốt)"
%% trang_thai: chinh-thuc
%% ghi_chu: "Kết luận gốc ý d) 'dùng số trung bình chính xác hơn' sai vì mẫu có giá trị bất thường; đúng là dùng trung vị (khớp DS Câu 5 cùng tài liệu). Tự duyệt theo yêu cầu giáo viên 22/10/2026."
\begin{ex}
Bác Dũng và bác Thu ghi lại số cuộc điện thoại mỗi người gọi mỗi ngày trong $10$ ngày được chọn ngẫu nhiên như sau:
\begin{center}
\begin{tabular}{|c|c|c|c|c|c|c|c|c|c|c|}\hline
Bác Dũng & $2$ & $7$ & $3$ & $6$ & $1$ & $4$ & $1$ & $4$ & $5$ & $1$\\\hline
Bác Thu & $1$ & $3$ & $1$ & $2$ & $3$ & $4$ & $1$ & $2$ & $20$ & $2$\\\hline
\end{tabular}
\end{center}
a) Hãy tìm số trung bình, các tứ phân vị và mốt của số cuộc điện thoại mỗi bác gọi.

b) Nếu so sánh theo số trung bình thì ai gọi nhiều cuộc điện thoại hơn?

c) Nếu so sánh theo trung vị thì ai gọi nhiều cuộc điện thoại hơn?

d) Nên dùng số trung bình hay trung vị để so sánh xem ai gọi nhiều cuộc điện thoại hơn mỗi ngày? Vì sao?
\loigiai{a) Bác Dũng: sắp xếp $1$; $1$; $1$; $2$; $3$; $4$; $4$; $5$; $6$; $7$; $\overline x=3{,}4$; $Q_1=1$; $Q_2=3{,}5$; $Q_3=5$; mốt $1$.

Bác Thu: sắp xếp $1$; $1$; $1$; $2$; $2$; $2$; $3$; $3$; $4$; $20$; $\overline x=3{,}9$; $Q_1=1$; $Q_2=2$; $Q_3=3$; mốt $1$ và $2$.

b) Theo số trung bình, bác Thu gọi nhiều hơn ($3{,}9>3{,}4$).

c) Theo trung vị, bác Dũng gọi nhiều hơn ($3{,}5>2$).

d) Nên dùng trung vị, vì số liệu của bác Thu có giá trị bất thường ($20$) làm số trung bình tăng lên, còn trung vị không bị ảnh hưởng.}
\end{ex}
